On the domain
$\Omega = \{ (\varphi,r) \mid 0\le \varphi < \pi \wedge 1<r<2 \}$
in polar coordinates,
two nonegative functions $f(r)$ and $g(r)$ are given.
The partial differential equation
\begin{equation}
\frac{\partial^2 u}{\partial r^2}
+
f(r)
\frac{\partial u}{\partial r}
+
g(r)
\frac{\partial^2 u}{\partial \varphi^2}
=
\lambda u
\label{40000037:pde}
\end{equation}
is to be solved with homogeneous boundary conditions.
Convert this equation into two ordinary differential equations with
appropriate boundary conditions and solve at least one of them.

\begin{loesung}
The domain is a ``rectangle'' in polar coordinates, so we use 
the separation ansatz $u(\varphi,r) = R(r) \Phi(\varphi)$ and
substitute it into the differential equation \eqref{40000037:pde}
giving
\begin{align*}
R''(r)\Phi(\varphi)
+
f(r) R'(r)\Phi(\varphi)
+
g(r) R(r)\Phi''(\varphi)
&=
\lambda R(r) \Phi(\varphi).
\intertext{Dividing by $R(r)\Phi(\varphi)$ gives}
\frac{R''(r)}{R(r)}
+
f(r) \frac{R'(r)}{R(r)}
+
g(r) \frac{\Phi''(\varphi)}{\Phi(\varphi)}
&=
\lambda.
\end{align*}
To separate variables, we need to devide by $g(r)$, which does not
cause any problems because $g(r)>0$.
Moving $\varphi$-terms to the right hand side gives the equation
\begin{equation}
\frac{1}{g(r)}
\biggl(
\frac{R''(r)}{R(r)}
+
f(r)
\frac{R'(r)}{R(r)}
-
\lambda
\biggr)
=
-\frac{\Phi''(\varphi)}{\Phi(\varphi)}.
\label{40000037:sepeq}
\end{equation}
The left hand side only depends on $r$, the right hand side only
on $\varphi$.
Therefore both sides are constant, we call this constant $\mu^2$
and separate
\eqref{40000037:sepeq}
and get
\begin{align*}
\frac{1}{g(r)}
\biggl(
\frac{R''(r)}{R(r)}
+
f(r)
\frac{R'(r)}{R(r)}
-
\lambda
\biggr)
&=
\mu^2
&&\text{and}
&
\mu^2
&=
-\frac{\Phi''(\varphi)}{\Phi(\varphi)}
\intertext{or}
R''(r)
+
f(r)R'(r)
-
(
\lambda
+
\mu^2 g(r)
)
R(r)
&=
0
&&\text{and}
&
\Phi''(\varphi)&=-\mu^2\Phi(\varphi).
\end{align*}
Appropriate boundary conditions are
\[
\begin{aligned}
R(1) &= 0 \\
R(2) &= 0 
\end{aligned}
\qquad\text{and}\qquad
\begin{aligned}
 \Phi(0)   &= 0 \\
 \Phi(\pi) &= 0.
\end{aligned}
\]
The equation for $R(r)$ can obviously not be solved without additional
knowledge of the functions $f(r)$ und $g(r)$, but the equation for
$\Phi(\varphi)$ is a harmonic oscillator equation that can be solved
by functions of the form
\[
\Phi(\varphi)
=
A\cos\mu \varphi + B\sin\mu\varphi
\]
with frequency $\mu$.
The boundary conditions imply
\[
\renewcommand{\arraycolsep}{1.5pt}
\begin{array}{lrccclcrclcrcl}
\text{at $\varphi=0$:}\quad&
0 &=& \Phi(0)   &=& A
&\quad\Rightarrow\quad&
A&=&0
\\
\text{at $\varphi=\pi$:}\quad&
0 &=& \Phi(\pi) &=& A\cos\mu\pi + B \sin\mu\pi
&\quad\Rightarrow\quad&
B\sin\mu\pi &=& 0
&\quad\Rightarrow\quad&
\mu\pi &=& k\pi,
\end{array}
\]
where $k\in\mathbb{Z}$.
Therefore the possible values for $\mu$ are $\mu=k\in\mathbb{Z}$
with solutions $\Phi_k(\varphi)=\sin k\varphi$.
\end{loesung}

\begin{diskussion}
To find a complete solution, the equation
\[
R''(r) + f(r)R'(r) - (\lambda+k^2g(r)) R(r) = 0
\]
needs to be solved with boundary conditions $R(1)=R(2)=0$.
For each $k$, this will only be possible for a discrete set of
values $\lambda_{lk}$ for the parameter $\lambda$, with
$l\in\mathbb{N}$, and will give solution functions $R_{lk}(r)$.
The eigenfunctions of the differential operator in
\eqref{40000037:pde}
for eigenvalue $\lambda_{lk}$ are thus
\[
u_{lk}(r,\varphi)
=
R_{lk}(r)\Phi_k(\varphi)
=
R_{lk}(r)\sin k\varphi.
\]
\end{diskussion}

\begin{bewertung}
Product ansatz ({\bf A}) 1 point,
separation of variables ({\bf S}) 1 point,
two separated differential equations ({\bf D}) 1 point,
boundary conditions ({\bf B}) 1 point,
solutions of the equation for $\Phi(\varphi)$ ({\bf L}) 1 point,
restriction of the value of $\mu$ from boundary conditions ({\bf R}) 1 point.
\end{bewertung}

